%
% Documento: Conclusão
\chapter{Conclusión}\label{chap:conclusion}

Se lograron realizar librerías basadas en Espressif IDF que pueden ser implementadas en módulos del esp32 como el esp32-wroom-32, esp32-wroom-32D y el esp32-s3-wroom-1. Estas mismas pueden manejar, a través del driver DM542T, motores paso a paso de 2 y 4 fases que no sobrepasen una corriente de 4.5A pico; cabe destacar que existe un límite máximo y mínimo de RPM del motor según la configuración de interruptores DIP que se le asigne al driver.Además, gracias a la librería encoder es posible obtener los grados efectivos del mismo para la aplicación que se necesite.

Debido a las diversas aplicaciones de los motores a paso es posible que estas librerías sean utilizadas en proyectos de robótica, CNC, impresoras 3D, entre otros. A pesar de esto faltó desarrollar considerablemente la aplicación de control en lazo cerrado para un sistema compuesto de péndulo invertido, el cual se dejó como trabajo futuro.

Por otra parte, el módulo desarrollado aunque funcional y estable, puede ser optimizado en cuanto a la gestión de memoria y el consumo energético, ya que el mismo no fue un foco principal durante el desarrollo del proyecto; e incluso relacionarlo a aplicaciones IoT.

Gracias al trabajo realizado, se logró adquirir experiencia en el desarrollo de software embebido para microcontroladores de la familia ESP32, así como en la implementación de controladores para motores paso a paso y la integración de interfaces de usuario mediante pantallas LCD y consolas serie. Esto permitió fortalecer habilidades en programación concurrente utilizando tareas y en el diseño de sistemas modulares y escalables.